% Part III, mechanising low-level program calculation: ~10 min
\partframe{Part III}{Mechanising Low-Level Program Calculation}

\begin{frame}[vertical-alignment = top]
  \frametitle{Two Ways Out of Rocq}
  \vspace*{10mm}
  \begin{center}
    \begin{tikzpicture}[alt={A specification is calculated by laws into
        a high-level program. Today, the high-level program is extracted to
        OCaml, relying on garbage-collected lists and trees, closures and
        boxed integers. The goal: more laws take the high-level program down
        to a low-level form, which is printed directly as low-level code
        with flat buffers, loops, machine integers, no allocation, and
        parallel hardware. The specification, the high-level program and the
        low-level form sit in a box labelled verified in Rocq.},
      stage/.style = {box, minimum height = 10mm, text width = 24mm,
        align = center},
      lbl/.style = {font = \bfseries},
      rocq/.style = {draw = accent, fill = accent!6, rounded corners = 6pt,
        thick, inner xsep = 3mm, inner ysep = 6mm}]
      \node[stage] (s) at (0, 0) {specification};
      \node[stage, draw = good, fill = good!8] (h) at (4.6, 0)
        {high-level program};
      \node[stage, right = 20mm of h] (o) {OCaml};
      \node[lbl, above = 1mm of o] {today};
      \draw[->, thick] (s) -- node[above, note] {laws} (h);
      \draw[->, thick] (h) -- node[pos = .6, above, note] {extract} (o);
      \node[note, right = 3mm of o, align = left, text width = 28mm]
        {GC'd lists, trees,\\ closures,\\ boxed integers};
      \uncover<2->{
        \node[stage, draw = good, fill = good!8] (lf) at (4.6, -2.6)
          {low-level form};
        \node[hot, minimum height = 10mm, text width = 24mm,
          align = center, right = 20mm of lf] (l) {\textbf{low-level code}};
        \node[lbl, above = 1mm of l] {goal};
        \draw[->, thick] (h) -- node[right, note] {laws} (lf);
        \draw[->, thick] (lf) -- node[pos = .6, above, note] {print} (l);
        \node[note, right = 3mm of l, align = left, text width = 28mm]
          {flat buffers, loops,\\ machine integers,\\
          no allocation,\\ parallel hardware};
      }
      \begin{scope}[on background layer]
        \only<1>{\node[rocq, fit = (s) (h)] (r) {};}
        \only<2->{\node[rocq, fit = (s) (h) (lf)] (r) {};}
        \node[note, text = accent, anchor = south west, font = \small\bfseries]
          at (r.south west) {\ verified in Rocq};
      \end{scope}
    \end{tikzpicture}
  \end{center}
  \ongoingtag
  \speakernote{Box = verified in Rocq \\
    Today: exit through OCaml (GC, closures, boxed integers) \\
    Goal: more laws down to a low-level form, printed directly \\
    Layout and scheduling become calculation steps \\
    Two lines of work show what such laws look like}
\end{frame}

\begin{frame}
  \frametitle{Trees Without Pointers}
  \begin{center}
    \begin{tikzpicture}[alt={Left: the tree Node of Node of Leaf 1 and
        Leaf 2, and Leaf 3, as heap cells linked by pointers. Right: the
        same tree serialised in preorder into one buffer: N, N, L, 1, L, 2,
        L, 3.},
      level distance = 10mm,
      level 1/.style = {sibling distance = 18mm},
      level 2/.style = {sibling distance = 12mm},
      every node/.style = {box, font = \small},
      cell/.style = {box, minimum width = 7mm, rounded corners = 0pt,
        font = \small\ttfamily}]
      \node {N}
        child {node {N} child {node {L 1}} child {node {L 2}}}
        child {node {L 3}};
      \node[draw = none, fill = none, note] at (0, 0.8) {pointers};
      \draw[->, very thick] (2.4, -1) -- node[above, draw = none,
        fill = none, note] {serialise} (4.0, -1);
      \begin{scope}[xshift = 45mm, yshift = -10mm, node distance = 0mm]
        \node[cell] (c0) {N};
        \foreach \v [count = \i from 1, remember = \i as \j (initially 0)]
            in {N, L, 1, L, 2, L, 3}
          \node[cell, right = of c\j] (c\i) {\v};
        \node[draw = none, fill = none, note, above = 2mm of c3]
          {one buffer, in preorder};
      \end{scope}
    \end{tikzpicture}
  \end{center}
  \begin{itemize}[<+->]
    \item Gibbon: compile traversals to run on the \textbf{buffer} itself.
    \item LoCal: types for locations and regions; parallel traversals;
      layouts chosen by a \textbf{cost model} (Marmoset).
  \end{itemize}
  \slideref{M. Vollmer et al. \textbf{Compiling Tree Transforms to Operate on Packed Representations}. ECOOP 2017. \rfsep
    M. Vollmer et al. \textbf{LoCal: A Language for Programs Operating on Serialized Data}. PLDI 2019. \rfsep
    C. Koparkar et al. \textbf{Efficient Tree-Traversals: Reconciling Parallelism and Dense Data Representations}. ICFP 2021. \rfsep
    V. Singhal et al. \textbf{Optimizing Layout of Recursive Datatypes with Marmoset}. ECOOP 2024.}
  \speakernote{Michael Vollmer's line of work (Gibbon compiler) \\
    Pointer trees: allocation, poor locality \\
    Serialised: one buffer, read in order \\
    LoCal: locations as types; ICFP'21: parallelism on dense data \\
    Marmoset: a cost model picks the field order}
\end{frame}

\begin{frame*}
  \frametitle{A Traversal over a Buffer}
  \mainblock[before upper = {\fvset{fontsize = \footnotesize}}]{packed-tree.py}
  \mainblock[before upper = {\fvset{fontsize = \footnotesize}}]{packed-sum.c}
  \simplified
  \uncover<2->{The right subtree starts \textbf{where the left one
    ends}.}
  \speakernote{Same function, on the buffer (hand-written sketch, not
    Gibbon's output) \\
    Left child: right after the tag \\
    Right child: only known after the left is traversed (end witnesses,
    LoCal) \\
    Order of traversal = order of layout}
\end{frame*}

\begin{frame*}
  \frametitle{Packed Traversals, Read as Hylomorphisms}
  \mainblock{packed-hylo.py}
  \begin{itemize}[<+->]
    \item Reading the buffer: \code{divide}. Summing: \code{conquer}.
      Together: \textbf{no tree}.
    \item Writing a buffer: a \code{conquer} step that \textbf{writes},
      in order.
    \item The catch: the position in the buffer flows from one child to
      the \textbf{next}.
  \end{itemize}
  \ongoingtag
  \speakernote{Our reading, not Gibbon's formulation \\
    Spec: unpack the buffer, fold the tree \\
    Program: traverse the buffer; the tree never exists \\
    Writes: an associative combiner (Part II) \\
    Cursor threading: state between siblings; laws beyond plain hylos}
\end{frame*}

\begin{frame*}
  \frametitle{Nested Loops, Divided and Conquered}
  \mainblock{nested-loop.c}
  \begin{itemize}[<+->]
    \item \textbf{Memoryless}: a row's result ignores the rows before:
      $s_{i+1} = s_i \oplus g(\mathit{row}_i)$.
    \item \textbf{Thm 4.7}: the nest splits like a homomorphism $\iff$
      memoryless, and the outer loop over row results does.
    \item \textbf{Thm 5.3}: extra state found for the outer loop works for
      the whole nest.
  \end{itemize}
  \slideref{A. Farzan, V. Nicolet. \textbf{Modular Divide-and-Conquer Parallelization of Nested Loops}. PLDI 2019.}
  \speakernote{Read-only multidimensional data; a loop nest computes a
    function \\
    Summarise the inner loop, then parallelise the outer one \\
    Sound and complete: nothing lost by going modular \\
    ParSynt: joins by SyGuS; checked in Rosette (bounded), most also in
    Dafny \\
    About 9--13x on 16 threads}
\end{frame*}

\begin{frame}
  \frametitle{The Same Story, with Hylomorphisms}
  \begin{center}
    \renewcommand\arraystretch{1.35}
    \begin{tabular}{ll}
      \textbf{Farzan, Nicolet} & \textbf{hylomorphisms} \\ \hline
      the inner loop, per row & an inner divide and conquer, per row \\
      memoryless & the outer \code{conquer}: a fold after a map \\
      homomorphism & an associative combiner \\
      lifting & extra state \\
      divide and conquer & \textbf{a new divide step}: split in $k$
        pieces \\
    \end{tabular}
  \end{center}
  \uncover<2->{Theorem 4.7 is the law our planned extension needs:
    \textbf{nested} hylos, parallelised one level at a time.}
  \ongoingtag
  \speakernote{Each notion has a recursion-scheme reading \\
    Homomorphism = associative combiner (monoid algebra, Part II) \\
    Divide and conquer = new divide step (Part I, same combiner) \\
    Their theorem justifies parallelising nested hylos level by level \\
    In Rocq: a law, proved once; no per-program bounded check}
\end{frame}

\begin{frame}
  \frametitle{Why a Mechanised Framework?}
  \begin{center}
    \small
    \renewcommand\arraystretch{1.35}
    \begin{tabular}{@{}llll@{}}
      & \textbf{the laws} & \textbf{each program} & \textbf{trusted} \\
      \hline
      ICFP 2016 & axioms & type checking & laws, code generator \\
      Gibbon & compiler passes & --- & the compiler \\
      ParSynt & paper proofs & join: bounded check, Dafny &
        translation, code generator \\
      \textbf{Rocq} & \textbf{proved once} & \textbf{every step checked} &
        \textbf{extraction or printer} \\
    \end{tabular}
  \end{center}
  \begin{itemize}[<+->]
    \item Same laws, one small trusted base.
    \item Clean extraction: the code that \textbf{runs} is the code that
      was \textbf{proved}.
  \end{itemize}
  \speakernote{ICFP'16: laws assumed \\
    Gibbon: correct by careful engineering, unverified \\
    ParSynt: theorems on paper; each join checked; loop translation not
    verified \\
    Ours: laws are Rocq theorems; derivation is a certificate \\
    Without clean extraction, this is only a proof about a model}
\end{frame}

\begin{frame}
  \frametitle{Laws That Change Representation}
  \begin{center}
    \begin{tikzpicture}[alt={Three kinds of rewrite. A tree of
        intermediate results becomes a flat buffer. Recursion becomes a
        loop. One traversal becomes independent pieces that can run at the
        same time.},
      from/.style = {box, text width = 32mm, align = center,
        minimum height = 10mm},
      to/.style = {box, draw = good, fill = good!8, text width = 32mm,
        align = center, minimum height = 10mm}]
      \foreach \y/\f/\t/\w/\o in {
          0/{intermediate tree}/{flat buffer}/{Gibbon}/1,
          -1.6/{recursion}/{loop}/{knapsack}/2,
          -3.2/{one traversal}/{independent pieces, in parallel}/{PLDI 2019}/3} {
        \uncover<\o->{
          \node[from] (f\o) at (0, \y) {\f};
          \node[to, right = 30mm of f\o] (t\o) {\t};
          \draw[->, thick] (f\o) -- node[above, note] {$=$, proved once}
            (t\o);
          \node[note, right = 3mm of t\o, text = gray] {\w};
        }
      }
    \end{tikzpicture}
  \end{center}
  \uncover<4->{Today's laws change the \emph{algorithm}; these change
    \emph{where data lives} and \emph{how work is scheduled}.}
  \ongoingtag
  \speakernote{Tree to buffer: what Gibbon does, as a law \\
    Recursion to loop: dyna\_loop already does this for knapsack \\
    Independent pieces: Farzan and Nicolet, as a law \\
    All proved once, used as rewrites}
\end{frame}

\begin{frame}
  \frametitle{Untrusted Search, Trusted Checking}
  \begin{center}
    \begin{tikzpicture}[alt={A specification and a cost model go into a
        search, which may use heuristics, equality saturation or a learned
        model, and is not trusted. It proposes a chain of laws. Rocq checks
        the chain, and the result is printed as low-level code. A wrong
        proposal is rejected.},
      stage/.style = {minimum height = 12mm, text width = 24mm,
        align = center}]
      \node[box, stage] (s) {specification\\ + cost model};
      \node[box, stage, right = 8mm of s, draw = bad, fill = bad!6] (q)
        {\textbf{search}\\ \small heuristics, e-graphs, ML};
      \node[box, stage, right = 22mm of q, draw = good, fill = good!8] (r)
        {\textbf{Rocq checks}\\ \small each step};
      \node[hot, stage, right = 8mm of r] (o) {low-level code};
      \draw[->, thick] (s) -- (q);
      \draw[->, thick] (q) -- node[above, note, align = center]
        {chain of\\ laws} (r);
      \draw[->, thick] (r) -- (o);
      \node[note, below = 2mm of q, text = bad] {untrusted};
      \node[note, below = 2mm of r, text = good] {trusted, small};
      % the same path, invisible, so the picture keeps its size
      \path (r.north) to[bend right = 30]
        node[above, note, opacity = 0] {rejected: try again} (q.north);
      \uncover<2->{
        \draw[->, thick, bad, dashed] (r.north) to[bend right = 30]
          node[above, note, text = bad] {rejected: try again} (q.north);
      }
    \end{tikzpicture}
  \end{center}
  \begin{itemize}
    \item The derivation is a \textbf{certificate}.
    \item<2-> The cost model of ICFP 2016, Marmoset's layouts, ParSynt's
      synthesis: all \textbf{search}.
  \end{itemize}
  \ongoingtag
  \speakernote{Search: heuristics, e-graphs, or an ML model \\
    Rocq checks the chain of laws; wrong proposals rejected \\
    Everything clever so far is search: keep it, check its answer}
\end{frame}

\begin{frame}
  \frametitle{Where We Are}
  \begin{center}
    \begin{tikzpicture}[alt={An index range from 0 to n is split into four
        pieces of different lengths. The pieces cover the range exactly,
        with no gaps and no overlaps; this is proved.},
      piece/.style = {draw = good, fill = good!8, minimum height = 7mm,
        inner sep = 0pt}]
      \node[piece, minimum width = 30mm, anchor = west] (p0) at (0, 0) {};
      \node[piece, minimum width = 18mm, right = 0mm of p0] (p1) {};
      \node[piece, minimum width = 24mm, right = 0mm of p1] (p2) {};
      \node[piece, minimum width = 26mm, right = 0mm of p2] (p3) {};
      \node[note, below = 1mm of p0.south west] {$0$};
      \node[note, below = 1mm of p3.south east] {$n$};
      \node[note, right = 3mm of p3] {exactly covered, no overlaps};
    \end{tikzpicture}
  \end{center}
  \begin{columns}[t]
    \begin{column}{0.5\textwidth}
      \textbf{\textcolor{good}{In Rocq, axiom-free:}}
      \begin{itemize}
        \item Splitting index ranges $k$ ways
        \item Machine integers, no overflow (proved)
        \item Buffers, compared by their reads
        \item Tree reductions over ranges
        \item Writes to a buffer, driven by any split
      \end{itemize}
    \end{column}
    \begin{column}{0.45\textwidth}
      \textbf{\textcolor{gray}{Not yet:}}
      \begin{itemize}
        \item Laws that add or remove a buffer
        \item Each index written exactly once
        \item Nested ranges
        \item The low-level target and printer
        \item The cost model and search
      \end{itemize}
    \end{column}
  \end{columns}
  \ongoingtag
  \speakernote{Early stage; be honest \\
    Done: exact-cover splits, machine integers, buffers, tree reductions,
    writes; materialising into a buffer defined, no laws yet \\
    Not yet: buffer laws, exactly-once, nesting, target, cost model}
\end{frame}
